\section{Independence tests}
\label{sec:independence-tests}

\subsection{Null hypothesis and weight regime}
\label{subsec:test-framework}

All tests in this section concern the hypothesis
$H_0\colon \mathbb{P}_{XY}=\mathbb{P}_X\otimes\mathbb{P}_Y$ of independence.
The finite-sample identities of the preceding sections are conditional on the
observed weights and need no stochastic model for them. Null distributions do
depend on such a model, and we distinguish two settings.

In the first setting, the observations are independent and identically
distributed and the weights are exogenous: deterministic, or random but
independent of all observations. We condition on the weights, which then form
a triangular array, and write $s_{r,n} = \sum_{i=1}^n p_i^r$, so that
$n_{\mathrm{eff}} = s_{2,n}^{-1}$. The limits below require
\begin{align}
    s_{2,n}
    &\longrightarrow0,
    &
    \frac{\max_{1\leq i\leq n} p_i}{\sqrt{s_{2,n}}}
    &\longrightarrow0.
    \label{eq:weight-diffuseness}
\end{align}
The first condition says that the effective sample size grows. The second is a
Lindeberg-type diffuseness condition: no single case may contribute a
non-negligible fraction of the standard deviation of a weighted sum. Uniform
weights satisfy it, whereas a sequence dominated by a fixed number of
observations does not. It implies, for example,
$s_{4,n}/s_{2,n}^2\leq(\max_ip_i)^2/s_{2,n}\to0$. This is the setting of the
Pearson, Spearman, Kendall, Blomqvist, and Hoeffding tests.

In the second setting, the weights are functions of the predictor only. Under
$H_0$ the responses remain conditionally independent and identically
distributed given the predictors and weights, which is what the weighted
Chatterjee test in Section~\ref{subsec:chatterjee-test} uses. The distinction
matters for the symmetric statistics: if $p_i$ is proportional to a function
of $X_i$, reweighting preserves the null factorization, but a squared-weight
law enters the sampling variance and Kish scaling need not give a pivotal
statistic. Weights that use the responses, or both variables jointly, are
covered by neither argument.

Whenever the weights are fixed or depend only on the predictors, a
finite-sample alternative is available. Conditional on $(X_i,p_i)_{i=1}^n$
and on the unordered responses, the assignment of responses to predictors is
uniform under $H_0$, so recomputing a statistic over random response
permutations gives an exact conditional test, also for concentrated weights or
ties. The analytic results below avoid this repeated computation, but their
conditions must then be respected. In particular, ties caused by atoms of the
population distribution are distinct from accidental equalities in
finite-precision data, and results that require continuous marginals are
labeled as such.

For a statistic $Z$ with a standard normal null limit, the upper-, lower-, and
two-sided p-values are $1-\Phi(Z)$, $\Phi(Z)$, and $2\Phi(-|Z|)$. Hoeffding's
test is an upper-tail test against nondirectional dependence; calling it
two-sided means only that increasing, decreasing, and nonmonotone alternatives
all make the statistic large.

\subsection{Pearson, Spearman, Kendall, and Blomqvist tests}
\label{subsec:standard-tests}

The first-order null theory of the four standard coefficients is classical.
We record it to collect all tests in one place and to distinguish the limiting
statistics from the finite-sample transformations used by the software.

\begin{proposition}[First-order weighted null limits]
\label{prop:standard-null-limits}
Assume \eqref{eq:weight-diffuseness}, exogenous weights, $H_0$, and continuous,
nondegenerate marginals. For Pearson's coefficient, also assume positive
finite marginal variances and moments sufficient for the weighted Lindeberg
condition; for Blomqvist's coefficient, assume unique regular marginal
medians. Then
\begin{align*}
  \sqrt{n_{\mathrm{eff}}}\,\widehat\rho_{P,w}
  &\mathrel{\Longrightarrow} N(0,1),
  &
  \sqrt{n_{\mathrm{eff}}}\,\widehat\rho_{S,w}
  &\mathrel{\Longrightarrow} N(0,1),\\
  \frac{3}{2}\sqrt{n_{\mathrm{eff}}}\,\widehat\tau_w
  &\mathrel{\Longrightarrow} N(0,1),
  &
  \sqrt{n_{\mathrm{eff}}}\,\widehat\beta_w
  &\mathrel{\Longrightarrow} N(0,1).
\end{align*}
\end{proposition}

The proof is the weighted triangular-array central limit theorem applied to an
asymptotic linear representation
$\widehat\theta_w=\sum_i p_i\psi(Z_i)+o_p(\sqrt{s_{2,n}})$, where $\psi$ is
the first-order null score of the coefficient. With $X$ and $Y$ standardized
for Pearson's coefficient and $U=F_X(X)$, $V=F_Y(Y)$ for the rank
coefficients, the scores are
\begin{align*}
  \psi_P(X,Y)
  &= \frac{X-\mathbb{E}X}{\sqrt{\operatorname{Var}(X)}}
     \frac{Y-\mathbb{E}Y}{\sqrt{\operatorname{Var}(Y)}},
  &
  \psi_S(U,V)
  &= 12(U-1/2)(V-1/2),\\
  \psi_K(U,V)
  &= 2(2U-1)(2V-1),
  &
  \psi_B(U,V)
  &= \{2\mathbf{1}(U\leq1/2)-1\}
     \{2\mathbf{1}(V\leq1/2)-1\},
\end{align*}
and under independence $\psi_P$, $\psi_S$, and $\psi_B$ have variance one
while $\psi_K$ has variance $4/9$, which explains the four constants. For
Spearman's coefficient the score is the influence function of its copula
functional; for Kendall's coefficient it is twice the first projection of the
concordance kernel; for Blomqvist's coefficient, estimation of the medians has
no first-order effect under independence because the opposite marginal sign
has mean zero. The same calculation shows why the limits need more than
replacing $n$ by $n_{\mathrm{eff}}$ when the weights depend on the
observations.

The software applies finite-sample transformations rather than the raw
limits. With $N=n_{\mathrm{eff}}$, it uses
\begin{align*}
  Z_P^{\mathrm{app}}
  &= \operatorname{atanh}(\widehat\rho_{P,w})\sqrt{N-3},
  &
  Z_S^{\mathrm{app}}
  &= \operatorname{atanh}(\widehat\rho_{S,w})
     \sqrt{\frac{N-3}{1.06}},
  &
  Z_B^{\mathrm{app}}
  &= \operatorname{atanh}(\widehat\beta_w)\sqrt N.
\end{align*}
The Pearson transformation is Fisher's normal-theory approximation
\citep{fisher1915}, and the Blomqvist transformation is first-order equivalent
to the statistic in Proposition~\ref{prop:standard-null-limits}. The factor
$1.06$ is the finite-sample approximation of \citet{fieller1957}; since
$Z_S^{\mathrm{app}} = \sqrt N\,\widehat\rho_{S,w}/\sqrt{1.06}+o_p(1)$, its
normal calibration is slightly conservative in the limiting framework of the
proposition. The Pearson and Spearman formulas require $N>3$, and every
$\operatorname{atanh}$ transformation requires its coefficient to lie strictly
between $-1$ and $1$.

For Kendall's tau, the tie-free weighted numerator $S_w$ has an exact
conditional null variance, and with ties the classical permutation variance
depends on the marginal tie multiplicities \citep{kendall1945,kendall1947}.
The implementation rescales the masses to fractional pseudo-counts $q_i=Np_i$
and substitutes weighted pair and triplet masses for the integer tie counts in
that formula. The resulting studentized statistic $Z_K^{\mathrm{app}}$, given
in \eqref{eq:kendall-pseudocount-statistic}, has exact conditional variance
one for uniform masses, is first-order equivalent to
Proposition~\ref{prop:standard-null-limits} without ties, and has the correct
first-order limit when ties arise from atoms. For unequal masses it is a
finite-sample approximation, not an exact permutation identity.
Appendix~\ref{app:null-asymptotics} records the algebra.

\subsection{Weighted Hoeffding test}
\label{subsec:hoeffding-test}

Hoeffding's estimator is unbiased under $H_0$, so its null mean is zero. Its
first projection also vanishes, which makes the statistic second-order
degenerate and rules out a Gaussian limit at the $\sqrt{n_{\mathrm{eff}}}$
scale. \citet{hoeffding1948} derived the degenerate null limit of the
unweighted statistic, and \citet{blum1961} tabulated the equivalent
nonnegative quadratic form used for calibration. The following theorem shows
that the weighted statistic has the same limit after an explicit
effective-sample-size correction.

\begin{theorem}[Weighted Hoeffding null limit]
\label{thm:weighted-hoeffding-null}
For each $n$, let $(p_1,\ldots,p_n)$ be a deterministic mass vector with at
least five positive entries, and let $(X_i,Y_i)_{i=1}^n$ be independent and
identically distributed. Suppose the marginals are continuous and
\eqref{eq:weight-diffuseness} holds along the triangular array. Under $H_0$,
\begin{align}
  \frac{1-s_{2,n}}{s_{2,n}}\frac{\widehat D_w}{30}
  &\mathrel{\Longrightarrow}
  \sum_{j=1}^{\infty}\sum_{k=1}^{\infty}
  \frac{G_{jk}^2-1}{\pi^4j^2k^2},
  \label{eq:hoeffding-centered-limit}
\end{align}
where the $G_{jk}$ are independent standard normal variables. Equivalently,
\begin{align}
  \mathcal H_{n,w}
  &= \frac{\pi^4}{2}
     \left\{
       (n_{\mathrm{eff}}-1)\frac{\widehat D_w}{30}
       +\frac{1}{36}
     \right\}
  \mathrel{\Longrightarrow}
  \mathcal H
  = \frac12\sum_{j=1}^{\infty}\sum_{k=1}^{\infty}
    \frac{G_{jk}^2}{j^2k^2}.
  \label{eq:hoeffding-positive-limit}
\end{align}
Moreover,
\begin{align*}
  \mathbb{E}_0(\widehat D_w\mid p)
  &=0,
  &
  \operatorname{Var}_0(\widehat D_w\mid p)
  &= \frac{2}{9}
     \frac{s_{2,n}^2-s_{4,n}}{(1-s_{2,n})^2}
     +o(s_{2,n}^2).
\end{align*}
\end{theorem}

The limit $\mathcal H$ is the Blum--Kiefer--Rosenblatt law, so the weighted
test can use the classical tables. The effective sample size enters through
the factor $n_{\mathrm{eff}}-1$, which replaces $n-1$ in the unweighted
statistic, while the centering constant $1/36$ is the same as in the
unweighted case. The variance expansion shows that the weights enter the
second-order behavior only through $s_{2,n}$ and a correction $s_{4,n}$ that
is negligible under \eqref{eq:weight-diffuseness}. The same conclusions hold
conditionally, path by path, for random exogenous weights whose realized array
satisfies the assumptions.

The test rejects at asymptotic level $\alpha$ when $\mathcal H_{n,w}$ exceeds
the $(1-\alpha)$ quantile of $\mathcal H$. The limit law is continuous:
conditioning on every Gaussian variable except $G_{11}$ leaves a noncentral
squared normal, whose distribution is continuous. Under a fixed alternative
with $D>0$, the weighted law of large numbers for each distinct-index
component gives $\widehat D_w\to_p D$, so $\mathcal H_{n,w}\to_p\infty$ and
the test is consistent.

The proof is given in Appendix~\ref{app:hoeffding-null-proof} and proceeds in
three steps. First, we compute the second-order von Mises expansion of the
functional $D/30$ at the independence copula. Its first derivative vanishes,
and its mixed second derivative is the product kernel
$K(z,z')=k(u,u')k(v,v')$ with $k(a,b)=(a^2+b^2)/2-\max(a,b)+1/3$. Second, we
apply Hoeffding's canonical decomposition to each distinct-index component
$\widehat\theta_{m,w}$, with the coefficient of each index set replaced by
its normalized distinct-tuple mass. Two projection identities that link the
decomposition to $K$ give
\begin{align*}
  \frac{\widehat D_w}{30}
  &= \frac{2}{1-s_{2,n}}
     \sum_{1\leq i<l\leq n}p_ip_lK(Z_i,Z_l)+R_n,
\end{align*}
where $R_n$ collects the residual first projection, the coefficient error of
the quadratic term, and the projections of orders three to five;
orthogonality of canonical projections and \eqref{eq:weight-diffuseness} give
$\mathbb{E}_0(R_n^2\mid p)=o(s_{2,n}^2)$. Finally, we expand $K$ in its
eigenfunctions $2\cos(\pi ju)\cos(\pi kv)$ with eigenvalues
$\lambda_{jk}=1/(\pi^4j^2k^2)$, apply the weighted Lindeberg central limit
theorem to finitely many eigen-coordinates, and control the spectral tail
uniformly in $n$ by canonicality. The identities
$\sum_{j,k}\lambda_{jk}=1/36$ and $\sum_{j,k}\lambda_{jk}^2=1/8100$ supply
the centering constant and the tail bound.

The survival function of $\mathcal H$ has no elementary form. The software
linearly interpolates tabulated probabilities of the Blum--Kiefer--Rosenblatt
limit and uses an exponential formula outside the table; both steps are
numerical approximations. Its internal convention replaces the $1/36$ in
\eqref{eq:hoeffding-positive-limit} by $(1-s_{2,n})/36$; the difference is
$s_{2,n}/36$ before multiplication by $\pi^4/2$ and is asymptotically
negligible.

The theorem does not extend automatically to predictor-dependent weights
$p_i\propto w(X_i)$. The horizontal covariance kernel then involves
\begin{align*}
  \Gamma_{X,w}(x,x')
  &=\frac{
    \mathbb{E}\!\left[
      w(X)^2
      \{\mathbf{1}(X\leq x)-F_w(x)\}
      \{\mathbf{1}(X\leq x')-F_w(x')\}
    \right]
  }{\mathbb{E}\{w(X)^2\}},
\end{align*}
where $F_w$ is the distribution induced by first-power weights. The first- and
second-power weight laws generally differ, which changes both the eigenvalues
and the centering constant. Conditional response permutation is the safe
calibration in that setting.

\subsection{Weighted Chatterjee test}
\label{subsec:chatterjee-test}

Put the positive-mass observations in the predictor order of
Section~\ref{sec:weighted-chatterjee}. Under $H_0$ with a continuous
response, and conditional on that order and the masses, the variables
$U_i=F_Y(Y_i)$, $i=1,\ldots,m$, are independent and uniform, and the weak
weighted rank $R_{i,w}$ equals $F_w(U_i)$ for the weighted empirical
distribution $F_w(u)=\sum_{j=1}^m p_j\mathbf{1}(U_j\leq u)$. The path
statistic is therefore
$A_w=\sum_{i=1}^{m-1}p_i|F_w(U_{i+1})-F_w(U_i)|$, and the reported
coefficient is $\widehat\xi_w=1-A_w/C_w$. For inference we use the surrogate
\begin{align*}
  \widetilde\xi_w
  &=1-3A_w,
\end{align*}
which replaces the random weighted-rank denominator $C_w$ by its limit $1/3$
and admits a closed-form conditional projection variance. The following
theorem gives the exact conditional null mean of $A_w$, its first-order
variance, and the resulting normal test.

\begin{theorem}[Conditional weighted Chatterjee null theory]
\label{thm:weighted-chatterjee-null}
Suppose the observations are independent and identically distributed and the
response is continuous and independent of the predictor. The weights may be
fixed or functions only of the predictor. Conditional on the predictor order,
its response-independent tie breaking, and the $m\geq2$ positive masses, set
\begin{align*}
  s_{2,m}
  &=\sum_{i=1}^m p_i^2,
  &P_E
  &=\sum_{i=1}^{m-1}p_i,
  &
  S_E
  &=\sum_{i=1}^{m-1}p_i^2,\\
  S_A
  &=\sum_{i=1}^{m-2}p_ip_{i+1},
  &
  S_{EN}
  &=\sum_{i=1}^{m-1}p_i(p_i+p_{i+1}).
\end{align*}
Then the conditional mean of $A_w$ is exactly
\begin{align}
  \mu_{A,w}
  &=\sum_{i=1}^{m-1}p_i
    \left\{\frac13+\frac{p_i+p_{i+1}}6\right\}
   =\frac{P_E}{3}+\frac{S_{EN}}6.
  \label{eq:chatterjee-null-mean}
\end{align}
Define the projection variance
\begin{align}
  v_{A,w}^2
  &=\frac{S_E}{18}+\frac{S_A}{90}
    +\frac{P_E^2s_{2,m}}{45}
    -\frac{P_ES_{EN}}{45}.
  \label{eq:chatterjee-projection-variance}
\end{align}
If, along the realized triangular array, $s_{2,m}\to0$ and
$\max_ip_i/\sqrt{s_{2,m}}\to0$, then
\begin{align*}
  \operatorname{Var}_0(A_w\mid X,p)
  &=v_{A,w}^2+o(s_{2,m}),
  &
  Z_{\xi,w}
  &=\frac{
      \widetilde\xi_w-(1-3\mu_{A,w})
    }{3v_{A,w}}
  \mathrel{\Longrightarrow}N(0,1).
\end{align*}
Moreover, $\widehat\xi_w-\widetilde\xi_w=o_p(\sqrt{s_{2,m}})$, so the
reported coefficient and the inferential surrogate are first-order equivalent
under the same conditions.
\end{theorem}

The mean formula is exact for every $m\geq2$ and every mass vector. Its
correction $S_{EN}/6$ to the naive value $P_E/3$ is of order $s_{2,m}$, so it
is asymptotically negligible relative to the standard deviation, but it makes
the centering of the implemented statistic exact. The variance
\eqref{eq:chatterjee-projection-variance} is a projection variance: it agrees
with the exact conditional variance to first order, but it is not the exact
finite-permutation variance, which
Appendix~\ref{app:chatterjee-null-moments} gives in closed form. For uniform
masses, $v_{A,w}^2$ recovers the familiar limiting variance $2/(5m)$ of
$\widehat\xi$, whereas the exact finite variance of the ordinary coefficient
is $(m-2)(4m-7)/\{10(m-1)^2(m+1)\}$. Because the argument conditions on the
predictor, the weights may depend on the predictor, in contrast to
Theorem~\ref{thm:weighted-hoeffding-null}. The implemented test uses the
centering $1-3\mu_{A,w}$ and the standard error $3v_{A,w}$.

The proof is given in Appendix~\ref{app:chatterjee-null-moments}. The edge
increment $|F_w(U_{i+1})-F_w(U_i)|$ is the mass of the responses lying
between the two endpoints, half-open at the upper endpoint, and its exact
conditional mean follows from counting which of three independent uniforms is
the middle one. Projecting $A_w-\mu_{A,w}$ onto functions of single edges and
single nodes gives a one-dependent sum whose variance is $v_{A,w}^2$; the four
terms in \eqref{eq:chatterjee-projection-variance} are the edge variances,
the adjacent-edge covariances, the node variances, and the edge--node
covariances. The remainder has conditional second moment
$O(s_{2,m}^2+\max_ip_i\,s_{2,m}+\max_ip_i^2)=o(s_{2,m})$, and a central limit
theorem for one-dependent triangular arrays \citep{janson2021} gives the
normal limit. The equivalence of $\widehat\xi_w$ and $\widetilde\xi_w$
follows from the bound $|C_w-1/3|\leq2\max_ip_i$.

The natural alternative is upper-tailed. A two-sided normal p-value can be
reported, but a lower-tail rejection should not be interpreted as negative
association, because the population coefficient is nonnegative and
directional.

If the response is discrete or tied, the library retains
\eqref{eq:weighted-chatterjee} as the reported coefficient. For uniform
masses it uses the tie-adjusted asymptotic variance of \citet{chatterjee2021},
recorded in Appendix~\ref{app:chatterjee-null-moments}. Analytic inference
for a tied response with unequal masses is not covered by
Theorem~\ref{thm:weighted-chatterjee-null} and is deliberately left
unavailable; conditional permutation remains valid when the weights are fixed
or predictor-dependent.

\subsection{Summary of test statistics}

Table~\ref{tab:independence-tests} separates proved asymptotic pivots from
finite-sample transformations. Normal entries use the p-value convention of
Section~\ref{subsec:test-framework}; Hoeffding uses the upper tail of
$\mathcal H$.

\begin{table}[H]
  \centering
  \caption{Independence-test statistics and their calibration status.}
  \label{tab:independence-tests}
  \begin{tabular}{lll}
    \toprule
    Measure & Statistic & Calibration status \\
    \midrule
    Pearson
      & $\operatorname{atanh}(\widehat\rho_{P,w})\sqrt{N-3}$
      & Fisher approximation \\
    Spearman
      & $\operatorname{atanh}(\widehat\rho_{S,w})\sqrt{(N-3)/1.06}$
      & conservative approximation \\
    Kendall
      & $Z_K^{\mathrm{app}}$ in
        \eqref{eq:kendall-pseudocount-statistic}
      & exact tie variance only for uniform masses \\
    Blomqvist
      & $\operatorname{atanh}(\widehat\beta_w)\sqrt N$
      & first-order normal \\
    Hoeffding
      & $\mathcal H_{n,w}$ in \eqref{eq:hoeffding-positive-limit}
      & quadratic-series limit \\
    Chatterjee
      & $Z_{\xi,w}$ in Theorem~\ref{thm:weighted-chatterjee-null}
      & conditional first-order normal \\
    \bottomrule
  \end{tabular}
\end{table}
